%=============================================================================
% IEEE Conference-format report — Direct-Mapped Cache Controller (Phase I)
% Compile with: pdflatex IEEE_Report.tex   (needs the IEEEtran class)
%   - Overleaf: create a project, upload this file, set compiler to pdfLaTeX.
%   - Local (TeX Live): `pdflatex IEEE_Report.tex` twice (for references).
%=============================================================================
\documentclass[conference]{IEEEtran}
\IEEEoverridecommandlockouts

\usepackage{cite}
\usepackage{amsmath,amssymb}
\usepackage{graphicx}
\usepackage{booktabs}
\usepackage{array}
\usepackage{url}

\begin{document}

\title{Design and Analysis of a Direct-Mapped Cache\\ Controller in Verilog}


\maketitle

%-----------------------------------------------------------------------------
\begin{abstract}
The growing gap between processor and main-memory speeds---the ``memory
wall''---makes on-chip caching essential to modern computer performance. This
work presents the design, implementation, and analysis of a \emph{direct-mapped
cache controller} in Verilog HDL. The prototype implements a 1~KB cache with
16-byte blocks for a 32-bit byte-addressed memory, organised as 64 lines with a
22-bit tag, 6-bit index, and 4-bit offset. The design comprises eight RTL
modules---an address decoder, tag/data/valid storage arrays, a tag comparator,
a behavioural main-memory model, and a three-state controller FSM---integrated
into a single top-level and driven by a trace-based testbench. Three
representative access patterns (sequential, repeated, and random) are used to
measure hit ratio, miss ratio, and Average Memory Access Time (AMAT). The RTL
results are cross-validated against an independent Python reference model and
agree exactly: hit ratios of 75.00\%, 96.88\%, and 16.02\%, giving AMAT values
of 3.50, 1.31, and 9.40 cycles respectively for a 10-cycle miss penalty. The
results quantitatively demonstrate how temporal and spatial locality govern
cache effectiveness.
\end{abstract}

\begin{IEEEkeywords}
cache memory, direct-mapped cache, memory hierarchy, AMAT, locality of
reference, Verilog HDL, computer architecture.
\end{IEEEkeywords}

%-----------------------------------------------------------------------------
\section{Introduction}
Processor clock speeds have historically improved far faster than main-memory
latency. A modern CPU can execute hundreds of instructions in the time taken to
service a single DRAM access. Caches bridge this gap by exploiting
\emph{locality of reference}: programs tend to reuse recently accessed data
(temporal locality) and to access nearby addresses (spatial locality).

This project designs and analyses the simplest practical cache
organisation---the \emph{direct-mapped cache}, in which every memory block maps
to exactly one cache line. The objectives are to (i) implement a synthesizable
cache controller in Verilog, (ii) correctly derive and handle the
tag/index/offset address fields and the valid bit, (iii) model the hit and miss
(refill) paths with a finite-state machine, and (iv) measure hit/miss
statistics and AMAT across different memory-access traces to analyse
architectural behaviour.

%-----------------------------------------------------------------------------
\section{Literature Review}
The foundations of cache memory and the memory hierarchy are established in the
standard texts. Hennessy and Patterson~\cite{hennessy} formalise the AMAT model
and the ``three Cs'' classification of misses (compulsory, capacity, conflict),
which directly frames the analysis in this work. Patterson and
Hennessy~\cite{patterson} present the direct-mapped organisation and the
tag/index/offset address decomposition used here.

Smith's classic survey~\cite{smith} catalogues cache design trade-offs---block
size, associativity, and replacement---and shows how increasing block size
improves spatial-locality hit rates up to a point before pollution dominates.
The principle of locality itself traces to Denning's work on program behaviour
and working sets~\cite{denning}.

On the implementation side, Harris and Harris~\cite{harris} give a didactic HDL
treatment of memory systems and finite-state-machine controllers, which informs
the controller structure in this prototype. These sources motivate our chosen
configuration (a small, analysable direct-mapped cache) and our use of synthetic
traces that isolate temporal versus spatial locality, following the evaluation
methodology common in the architecture literature.

%-----------------------------------------------------------------------------
\section{Proposed Solution}

\subsection{Cache Organisation}
A 1~KB direct-mapped cache with 16-byte blocks is used. For a 32-bit byte
address the fields are shown below.

\begin{verbatim}
 31                      10 9     4 3    0
+--------------------------+-------+------+
|         TAG (22)         |INDEX 6|OFF 4 |
+--------------------------+-------+------+
\end{verbatim}

\begin{itemize}
  \item \textbf{Offset (4 bits)} $=\log_2(16)$: byte within a block; bits [3:2]
        select one of the four 32-bit words.
  \item \textbf{Index (6 bits)} $=\log_2(64)$: selects one of 64 lines.
  \item \textbf{Tag (22 bits)} $= 32-6-4$: stored per line and compared on each
        access.
\end{itemize}

\subsection{RTL Modules}
Table~\ref{tab:modules} lists the eight RTL modules.

\begin{table}[htbp]
\caption{RTL modules and their functions}
\label{tab:modules}
\centering
\begin{tabular}{@{}ll@{}}
\toprule
\textbf{Module} & \textbf{Function} \\
\midrule
\texttt{address\_decoder}  & splits address into tag/index/offset/word-select \\
\texttt{tag\_ram}          & $64 \times 22$-bit stored tags \\
\texttt{cache\_data\_ram}  & $64 \times 128$-bit data blocks \\
\texttt{valid\_ram}        & 64 valid bits, cleared on reset \\
\texttt{tag\_comparator}   & \texttt{hit = valid AND (tag\_in == tag\_st)} \\
\texttt{main\_memory}      & behavioural model, 10-cycle fetch latency \\
\texttt{cache\_controller} & three-state FSM (lookup and refill) \\
\texttt{cache\_top}        & integrates datapath, controller, memory \\
\bottomrule
\end{tabular}
\end{table}

\subsection{Controller FSM}
The controller is a Mealy FSM with three states:
\begin{itemize}
  \item \textbf{IDLE} --- the hit signal is combinational, so a hit is answered
        in the same cycle (\texttt{cpu\_ready}), while a miss asserts
        \texttt{mem\_req} and moves to ALLOCATE.
  \item \textbf{ALLOCATE} --- holds the memory request until the block arrives,
        then asserts the write-enable to load the data, tag, and valid bit, and
        moves to DONE.
  \item \textbf{DONE} --- returns the requested word to the CPU and goes back to
        IDLE.
\end{itemize}
A hit costs \textbf{1 cycle}; a miss costs \textbf{$1 + 10 = 11$ cycles}. The
10-cycle penalty is a tunable parameter representing main-memory latency.

\subsection{Verification Methodology}
A trace-driven testbench issues each address from a trace file, waits for the
response, and counts hit/miss strobes. An independent Python reference model
(\texttt{cache\_sim.py}) reproduces the same tag/index/offset logic and
cold-start behaviour; agreement between the two is the correctness criterion.

%-----------------------------------------------------------------------------
\section{Results and Discussion}
Three traces of 256 accesses each were run. RTL (Icarus Verilog) and the
reference model produced identical numbers, shown in Table~\ref{tab:results}.

\begin{table}[htbp]
\caption{Hit/miss statistics and AMAT (256 accesses, 10-cycle penalty)}
\label{tab:results}
\centering
\begin{tabular}{@{}lrrrrr@{}}
\toprule
\textbf{Trace} & \textbf{Hits} & \textbf{Miss} & \textbf{Hit\%} & \textbf{Miss\%} & \textbf{AMAT} \\
\midrule
Sequential & 192 & 64  & 75.00 & 25.00 & 3.500 \\
Repeated   & 248 & 8   & 96.88 & 3.12  & 1.312 \\
Random     & 41  & 215 & 16.02 & 83.98 & 9.398 \\
\bottomrule
\end{tabular}
\end{table}

The AMAT model is
\begin{equation}
\text{AMAT} = t_{hit} + m \times t_{penalty},
\end{equation}
with $t_{hit}=1$ cycle, $t_{penalty}=10$ cycles, and $m$ the measured miss rate.

\textbf{Discussion.} The \emph{sequential} trace strides one word at a time;
because each 16-byte block holds four words, exactly one access in four
misses---a direct measurement of \emph{spatial locality} (75\% hit). The
\emph{repeated} trace cycles over a small working set that fits in the cache, so
after the compulsory misses almost every access hits (96.88\%)---\emph{temporal
locality}. The \emph{random} trace spreads accesses over four times the cache
size, so conflict and capacity misses dominate (16.02\% hit), and AMAT
approaches the full miss penalty. AMAT thus ranges from 1.31 cycles (good
locality, memory latency largely hidden) to 9.40 cycles (poor locality, penalty
paid almost every access), exactly as the memory-hierarchy model predicts.

%-----------------------------------------------------------------------------
\section{Individual Contributions}
\begin{itemize}
  \item \textbf{Author 1} --- Architecture and address decoding: cache parameter
        definitions, the address-decoder module, architecture documentation, and
        repository setup.
  \item \textbf{Author 2} --- Cache storage arrays: the data RAM, tag RAM, and
        valid RAM, including reset handling.
  \item \textbf{Author 3} --- Controller FSM, behavioural main-memory model, tag
        comparator, top-level integration (\texttt{cache\_top}), and the
        trace-driven testbench.
  \item \textbf{Author 4} --- Trace generation, the Python reference model, AMAT
        analysis, and report preparation.
\end{itemize}
The project's version-control commit history reflects these contributions.

%-----------------------------------------------------------------------------
\section{Conclusion and Future Work}
A functionally correct direct-mapped cache controller was designed, implemented,
and verified in Verilog, and its behaviour was analysed across three access
patterns. The measured hit ratios and AMAT values match an independent model and
clearly demonstrate the roles of temporal and spatial locality. \emph{Phase II}
will extend this prototype with logic synthesis (area/timing results),
cache-size and block-size parameter sweeps, write-policy support, and a
direct-mapped versus set-associative comparison.

%-----------------------------------------------------------------------------
\begin{thebibliography}{1}
\bibitem{hennessy}
J.~L. Hennessy and D.~A. Patterson, \emph{Computer Architecture: A Quantitative
Approach}, 6th~ed. Cambridge, MA, USA: Morgan Kaufmann, 2019.

\bibitem{patterson}
D.~A. Patterson and J.~L. Hennessy, \emph{Computer Organization and Design: The
Hardware/Software Interface}, 5th~ed. Waltham, MA, USA: Morgan Kaufmann, 2014.

\bibitem{smith}
A.~J. Smith, ``Cache memories,'' \emph{ACM Computing Surveys}, vol.~14, no.~3,
pp.~473--530, Sep.~1982.

\bibitem{denning}
P.~J. Denning, ``The locality principle,'' \emph{Communications of the ACM},
vol.~48, no.~7, pp.~19--24, Jul.~2005.

\bibitem{harris}
D.~M. Harris and S.~L. Harris, \emph{Digital Design and Computer Architecture},
2nd~ed. Waltham, MA, USA: Morgan Kaufmann, 2012.
\end{thebibliography}

\end{document}
